% --- IMPORTS ---
\documentclass[leqno]{article}
\usepackage{graphicx}
\usepackage{dsfont}
\usepackage{mathtools}
\usepackage{amssymb}
\usepackage{amsmath}
\usepackage{amsthm}
\usepackage{float}
\usepackage{ragged2e}
\usepackage{tensor}
\usepackage{fancyhdr}
\usepackage{empheq}
\usepackage[most]{tcolorbox}
\usepackage{enumitem}
\usepackage{dirtytalk}
\usepackage{marginnote}
\usepackage{microtype}
\usepackage[
  a4paper,
  left=0.8in,
  right=1in,
  top=1in,
  bottom=0.5in,
  headheight=14pt,
  headsep=0.4in
]{geometry}
\setlength{\parindent}{0cm}
% --- TikZ Packages ---
\usepackage{pgfplots}
\pgfplotsset{compat=1.18}
\usepgfplotslibrary{fillbetween, polar}
\usetikzlibrary{calc, positioning}
\usetikzlibrary{angles, quotes}
\usetikzlibrary{arrows.meta}
\usetikzlibrary{decorations.pathreplacing, decorations.markings}
\usetikzlibrary{patterns}
\usetikzlibrary{intersections}
% --- URL Packages ---
\usepackage[colorlinks=true,linkcolor=red,citecolor=magenta,urlcolor=black]{hyperref}%
\urlstyle{same}
% --- END OF IMPORTS ---

% --- CUSTOM COMMANDS ---
\RenewDocumentCommand{\marks}{o m}{%
  \IfValueTF{#1}%
    {\marginnote{\raggedright\textbf{[#2]}}[#1]}%
    {\marginnote{\raggedright\textbf{[#2]}}}%
}
\newcommand{\vspacerule}[2]{%
  \par\vspace*{#1}%
  \noindent\hrulefill\par
  \vspace*{#2}%
}
\definecolor{scarlet}{RGB}{205, 40, 75}
\newcommand{\questionitem}{%
  \item\phantomsection
  \addcontentsline{toc}{section}{Question \number\value{enumi}}%
}
\renewcommand{\contentsname}{Questions}
% --- END OF CUSTOM COMMANDS ---

% --- HEADER CONFIG ---
\pagestyle{fancy}
\fancyhead{}
\fancyfoot{}
\renewcommand{\headrulewidth}{0.9pt}
\lhead{Roots of Polynomials}
\chead{}
\rhead{\href{https://danieljsmith.org}{danieljsmith.org}}
% --- END OF HEADER CONFIG ---

\begin{document}
\vspace*{20pt}
\tableofcontents
\newpage
\begin{enumerate}[
  label=\textbf{\arabic*.},
  leftmargin=*,
  topsep=0pt,
  partopsep=0pt,
  parsep=0pt
]

% ---- Question 1 ----
\questionitem

The cubic equation
\[
x^3 - 5x^2 + px - q = 0
\]
where $p$ and $q$ are positive real constants, has roots $\alpha$, $\beta$ and $\gamma$. \\

Given that
\[
(\alpha-\beta)^2 + (\beta-\gamma)^2 + (\gamma-\alpha)^2 = 14
\]
\begin{enumerate}[label=\textbf{(\alph*)}]
  \item show that $p=6$ \marks{3} \\
\end{enumerate}
Given that
\[
\frac{1}{\alpha} + \frac{1}{\beta} + \frac{1}{\gamma} = 3
\]
\begin{enumerate}[label=\textbf{(\alph*)}, resume]
  \item determine the value of $q$ \marks{3} \\
  \item Without solving the cubic equation, determine the value of $(\alpha+1)(\beta+1)(\gamma+1)$ \marks{4} \\
\end{enumerate}

\newpage

% ---- Question 2 ----
\questionitem

The quartic equation $x^4+bx^3+cx^2+dx+3=0$ has roots $\alpha$, $\beta$, $\gamma$, $\delta$. It is given that
\[
\alpha+\beta+\gamma+\delta=-2,\qquad (\alpha+1)^2+(\beta+1)^2+(\gamma+1)^2+(\delta+1)^2=6,\qquad \alpha^{-1}+\beta^{-1}+\gamma^{-1}+\delta^{-1}=2
\]
\begin{enumerate}[label=\textbf{(\alph*)}]
  \item Find the values of $b$, $c$ and $d$ \marks{6} \\
  \item Given also that
  \[
  (\alpha+1)^3+(\beta+1)^3+(\gamma+1)^3+(\delta+1)^3=20
  \]
  find the value of $\alpha^4+\beta^4+\gamma^4+\delta^4$ \marks{2} \\
\end{enumerate}

\newpage

% ---- Question 3 ----
\questionitem

The equation $x^3 - 4x^2 - x + 6 = 0$ has roots $\alpha$, $\beta$ and $\gamma$. \\

Determine the value of \[\frac{1}{\alpha+\beta} + \frac{1}{\beta+\gamma} + \frac{1}{\gamma+\alpha}\] \marks{4}

\newpage

% ---- Question 4 ----
\questionitem

The quartic equation $3x^4 - x^3 - 10x^2 + 8x + 4 = 0$ has roots $\alpha$, $\beta$, $\gamma$ and $\delta$. \\

Without solving the equation, find equations with integer coefficients whose roots are \\
\begin{enumerate}[label=\textbf{(\alph*)}]
  \item $\dfrac{1}{\alpha}$, $\dfrac{1}{\beta}$, $\dfrac{1}{\gamma}$ and $\dfrac{1}{\delta}$ \marks{6} \\
  \item $\dfrac{2}{\alpha - 1}$, $\dfrac{2}{\beta - 1}$, $\dfrac{2}{\gamma - 1}$ and $\dfrac{2}{\delta - 1}$ \marks{6} \\
\end{enumerate}

\newpage

% ---- Question 5 ----
\questionitem

The cubic equation $x^3 + x^2 - 2x - 1 = 0$ has roots $\alpha$, $\beta$ and $\gamma$. \\

\begin{enumerate}[label=\textbf{(\alph*)}]
  \item Find a cubic equation whose roots are $\frac{1}{\alpha+1}$, $\frac{1}{\beta+1}$ and $\frac{1}{\gamma+1}$ \marks{3} \\
  \item Find the value of \[\left(\frac{1}{\alpha+1}\right)^2 + \left(\frac{1}{\beta+1}\right)^2 + \left(\frac{1}{\gamma+1}\right)^2\] \marks[-20pt]{2} \\
  \item Find the value of \[\left(\frac{1}{\alpha+1}\right)^3 + \left(\frac{1}{\beta+1}\right)^3 + \left(\frac{1}{\gamma+1}\right)^3\] \marks{2} \\
\end{enumerate}

\newpage

% ---- Question 6 ----
\questionitem

The equation $x^3 - 14x^2 + 56x + c = 0$, where $c$ is a constant, has three roots in geometric progression. \\

\begin{enumerate}[label=\textbf{(\alph*)}]
  \item Determine the roots of the equation \marks{6} \\
  \item Find $c$ \marks{1} \\
\end{enumerate}

\newpage

% ---- Question 7 ----
\questionitem

The quadratic equation
\[
x^2 - kx + 1 = 0
\]
where $k$ is an integer, has roots $\alpha$ and $\beta$. \\

\begin{enumerate}[label=\textbf{(\alph*)}]
  \item Write down, in terms of $k$ where appropriate, the value of $\alpha + \beta$ and the value of $\alpha\beta$ \marks{2} \\
  \item Determine, in simplest form in terms of $k$, the value of
  \[
  (\alpha^2 + \beta) + (\beta^2 + \alpha)
  \]
  \marks[-20pt]{3}
  \item Determine a quadratic equation which has roots
  \[
  \alpha^2 + \beta \text{ and } \beta^2 + \alpha
  \]
  giving your answer in the form $px^2 + qx + r = 0$, where $p$, $q$ and $r$ are integer values in terms of $k$ \marks{4} \\
\end{enumerate}

\newpage

% ---- Question 8 ----
\questionitem

The roots of the equation
\[
x^3 + 5x^2 - 4x - 24 = 0
\]
are $\alpha$, $\beta$ and $\gamma$. \\

Without solving the equation,
\begin{enumerate}[label=\textbf{(\alph*)}]
  \item write down the value of each of
  \[
  \alpha + \beta + \gamma \qquad \alpha\beta + \beta\gamma + \gamma\alpha \qquad \alpha\beta\gamma
  \]
  \marks[-20pt]{1}
  \item Hence determine the value of
  \begin{enumerate}[label=(\roman*)]
    \item \[
    \frac{1}{\alpha\beta} + \frac{1}{\beta\gamma} + \frac{1}{\gamma\alpha}
    \]
    \marks[-20pt]{2}
    \item \[
    (\alpha + 2)(\beta + 2)(\gamma + 2)
    \]
    \marks[-20pt]{2}
    \item \[
    \alpha^2\beta^2 + \beta^2\gamma^2 + \gamma^2\alpha^2
    \]
    \marks[-20pt]{3}
  \end{enumerate}
\end{enumerate}

\newpage

% ---- Question 9 ----
\questionitem

The roots of the equation $3x^3 + px^2 - 12x - 8 = 0$ are $\alpha$, $\beta$ and $\gamma$. \\

\begin{enumerate}[label=\textbf{(\alph*)}]
  \item Given that $\alpha + \beta + \gamma = 2$, write down the value of $p$ \marks{1} \\
  \item Write down values for $\alpha\beta + \beta\gamma + \gamma\alpha$ and $\alpha\beta\gamma$ \marks{1} \\
  \item Hence find the value of $\left(1+\frac{1}{\alpha}\right)\left(1+\frac{1}{\beta}\right)\left(1+\frac{1}{\gamma}\right)$ \marks{3} \\
\end{enumerate}

\newpage

% ---- Question 10 ----
\questionitem

The roots of the equation $x^3 + 2x^2 - 5x - 3 = 0$ are $\alpha$, $\beta$ and $\gamma$. \\

\begin{enumerate}[label=\textbf{(\alph*)}]
  \item Find $\alpha^2\beta^2 + \beta^2\gamma^2 + \gamma^2\alpha^2$ \marks{4} \\
  \item Find an equation with integer coefficients whose roots are $\alpha^2$, $\beta^2$ and $\gamma^2$ \marks{4} \\
\end{enumerate}

\newpage

% ---- Question 11 ----
\questionitem

The graph of $y=x^3+3x^2+px-20$, where $p$ is a constant, crosses the $x$-axis at three points whose $x$-coordinates are equally spaced. \\

Find the roots of
\[
x^3+3x^2+px-20=0
\] \marks[-20pt]{6}

\newpage

% ---- Question 12 ----
\questionitem

The cubic equation $x^3-3x+1=0$ has roots $\alpha$, $\beta$ and $\gamma$. \\

\begin{enumerate}[label=\textbf{(\alph*)}]
  \item Find a cubic equation whose roots are $\dfrac{1}{\alpha-1}$, $\dfrac{1}{\beta-1}$, $\dfrac{1}{\gamma-1}$. \marks{3} \\
  \item Hence find the value of $\dfrac{1}{(\alpha-1)^2}+\dfrac{1}{(\beta-1)^2}+\dfrac{1}{(\gamma-1)^2}$. \marks{2} \\
  \item Find also the value of $\dfrac{1}{(\alpha-1)^5}+\dfrac{1}{(\beta-1)^5}+\dfrac{1}{(\gamma-1)^5}$. \marks{3} \\
\end{enumerate}

\newpage

% ---- Question 13 ----
\questionitem

The quartic equation $x^4-2x^2+2x+1=0$ has roots $\alpha$, $\beta$, $\gamma$, $\delta$. \\

\begin{enumerate}[label=\textbf{(\alph*)}]
  \item Find a quartic equation whose roots are $\alpha^2$, $\beta^2$, $\gamma^2$, $\delta^2$ and state the value of $\alpha^2+\beta^2+\gamma^2+\delta^2$. \marks{5} \\
  \item Find the value of $\alpha^6+\beta^6+\gamma^6+\delta^6$. \marks{3} \\
  \item Find the value of $\alpha^8+\beta^8+\gamma^8+\delta^8$. \marks{2} \\
\end{enumerate}

\newpage

% ---- Question 14 ----
\questionitem

The equation $z^4-2z^3-z^2-2z+1=0$ has roots $\alpha$, $\beta$, $\gamma$ and $\delta$. \\

\begin{enumerate}[label=\textbf{(\alph*)}]
  \item Show that a quartic equation whose roots are $\alpha+\dfrac{1}{\alpha}$, $\beta+\dfrac{1}{\beta}$, $\gamma+\dfrac{1}{\gamma}$ and $\delta+\dfrac{1}{\delta}$ is
  \[
  w^4-4w^3-2w^2+12w+9=0
  \]
  \marks[-20pt]{3}
  \item Hence determine the exact roots of the equation $z^4-2z^3-z^2-2z+1=0$. \marks{3} \\
\end{enumerate}

\newpage

% ---- Question 15 ----
\questionitem

The quartic equation $x^4-2x^3+3x^2-4x+1=0$ has roots $\alpha$, $\beta$, $\gamma$, $\delta$. \\

\begin{enumerate}[label=\textbf{(\alph*)}]
  \item Find the value of $\dfrac{1}{\alpha}+\dfrac{1}{\beta}+\dfrac{1}{\gamma}+\dfrac{1}{\delta}$. \marks{2} \\
  \item Find the value of $\dfrac{1}{\alpha^2}+\dfrac{1}{\beta^2}+\dfrac{1}{\gamma^2}+\dfrac{1}{\delta^2}$. \marks{2} \\
  \item Find the value of
  \[
  \left(\dfrac{1}{\alpha}+\dfrac{1}{\beta}+\dfrac{1}{\gamma}\right)^2+\left(\dfrac{1}{\beta}+\dfrac{1}{\gamma}+\dfrac{1}{\delta}\right)^2+\left(\dfrac{1}{\gamma}+\dfrac{1}{\delta}+\dfrac{1}{\alpha}\right)^2+\left(\dfrac{1}{\delta}+\dfrac{1}{\alpha}+\dfrac{1}{\beta}\right)^2
  \]
  \marks[-20pt]{5}
\end{enumerate}
\end{enumerate}
\end{document}
